% !PRINTED 131
\weberpage{131}
%<<CONT
und daraus endlich:
\weq{3}{e^{-\pi i \beta v v'}\vartheta_{11}(v',\omega') = \varepsilon^3\sqrt{\alpha+\beta\omega}\,\vartheta_{11}(v,\omega).}

Die Transformationsformeln der drei übrigen Funktionen $\vartheta_{00}, \vartheta_{01}, \vartheta_{10}$ sind verschieden in den sechs Klassen des § 36 und können aus den Formeln (5), § 36 in derselben Weise hergeleitet werden. Man kann die sechs Fälle aber auch in ein einziges Formelsystem zusammenfassen, das man aus (3) erhält, wenn man $v$ ersetzt durch
\[
v+\frac{\alpha+\beta\omega}{2}, \quad v+\frac{\gamma+\delta\omega}{2}, \quad v+\frac{\alpha+\gamma+(\beta+\delta)\omega}{2},
\]
also $v'$ durch
\[
v'+\frac{1}{2}, \quad v'+\frac{\omega'}{2}, \quad v' = \frac{1+\omega'}{2},
\]
% [sic] Printed equals sign in the third member; mathematically expected +.
und dann auf der rechten und linken Seite von (3) die Formeln (2), (3), (8) des § 21 anwendet. So kommt:
\weq{4}{
\begin{aligned}
&e^{-\pi i \beta v v'}\vartheta_{10}(v',\omega') \\
&= i^{\beta} e^{-\frac{\pi i \alpha\beta}{4}} \varepsilon^3 \sqrt{\alpha+\beta\omega}\,\vartheta_{1+\beta,\,1-\alpha}(v,\omega),
\end{aligned}
}
\weq{5}{
\begin{aligned}
&e^{-\pi i \beta v v'}\vartheta_{01}(v',\omega') \\
&= i^{\delta-1} e^{-\frac{\pi i \gamma\delta}{4}} \varepsilon^3 \sqrt{\alpha+\beta\omega}\,\vartheta_{1+\delta,\,1-\gamma}(v,\omega),
\end{aligned}
}
\weq{6}{
\begin{aligned}
&e^{-\pi i \beta v v'}\vartheta_{00}(v',\omega') \\
&= i^{\delta+\beta-\alpha\delta} e^{-\frac{\pi i}{4}(\alpha\beta+\gamma\delta)} \varepsilon^3 \sqrt{\alpha+\beta\omega}\,\vartheta_{1+\beta+\delta,\,1-\alpha-\gamma}(v,\omega).
\end{aligned}
}

Die vierten Potenzen dieser Funktionen lassen sich einfacher ausdrücken durch
\weq{7}{e^{-4\pi i \beta v v'}\vartheta^4_{g_1,g_2}(v',\omega') = (-1)^{g_2\alpha\beta+g_1\gamma\delta+\beta\gamma}(\alpha+\beta\omega)^2\,\vartheta^4_{g_1',g_2'}(v),}
% [sic] Printed without the second argument $\omega$; possibly an intentional abbreviation.
worin in den sechs Klassen die zu den Charakteristiken $(g_1,g_2)=(00),(01),(10)$ gehörigen Charakteristiken $(g_1',g_2')$ aus der letzten Kolumne der Tabelle in § 36 zu entnehmen sind.

Eine einfachere Transformationsformel erhält man aus den $\sigma$-Funktionen für das Produkt der drei $\vartheta$-Funktionen (4), (5), (6). Es ist nämlich, wenn man das Produkt der drei Funktionen § 37, (6) bildet:
\[
\sigma_{00}(u)\,\sigma_{01}(u)\,\sigma_{10}(u)
= e^{\frac{3\eta_1 u^2}{\omega_1}}\; \frac{\vartheta_{00}\!\left(\dfrac{u}{\omega_1}, \dfrac{\omega_2}{\omega_1}\right)\vartheta_{01}\!\left(\dfrac{u}{\omega_1}, \dfrac{\omega_2}{\omega_1}\right)\vartheta_{10}\!\left(\dfrac{u}{\omega_1}, \dfrac{\omega_2}{\omega_1}\right)}{\vartheta_{00}\,\vartheta_{01}\,\vartheta_{10}}
\]
%CONT>>
